\chapter*{ДҮГНЭЛТ}
\addcontentsline{toc}{chapter}{ДҮГНЭЛТ}
\markboth{ДҮГНЭЛТ}{ДҮГНЭЛТ}

Энэхүү судалгааны ажлын хүрээнд дадал үүсэх онолын механизмыг програм хангамжийн шаардлага, архитектур, өгөгдлийн загвар, тайлбарлагдахуйц дүрэм болгон хөрвүүлж, зан үйлд суурилсан дадал төлөвшүүлэх системийг хэрэгжүүлэв. Судалгааны зорилтуудын биелэлтийг дор нэгтгэв.

\begin{enumerate}
    \item Дадал төлөвшлийн мөчлөг, Фоггийн зан үйлийн загвар, итгүүлэн нөлөөлөх системийн зохиомж, SRBAI хэмжүүр болон дадал үүсэх хугацааны судалгааг тоймлон, системийн зохиомжид ашиглах зарчмуудыг тодорхойлсон (2-р бүлэг).
    \item Loop Habit Tracker, Streaks, Habitica зэрэг одоо ашиглагдаж буй програмуудын албан ёсны тайлбарт үндэслэн тэдгээрийн боломжийг харьцуулж, өдөөгч нөхцөл, автоматжилтын хэмжилт, сануулгын бууралтын хувьд дутуу байгаа хэсгийг тодорхойлсон (2-р бүлэг).
    \item Системийн үзэл баримтлал, функциональ болон функциональ бус шаардлага, хэрэглээний тохиолдлыг тодорхойлсон (3, 4-р бүлэг, хавсралт).
    \item Өдөөгч дохио, хариу үйлдэл, урамшуулал, эргэцүүллийг цагийн хүрээ ба газрын нөхцөл, доод зорилт, өдөрт нэг гүйцэтгэлийн бүртгэл, сануулгын шийдвэр, гүйцэтгэлийн эх үүсвэрийн ангилал, XP ба амжилтын тэмдэг зэрэг тодорхой дүрэм, өгөгдлийн бүтэцтэй уялдуулан загварчилсан (3--5-р бүлэг).
    \item Дадлын хүчийг SRBAI автоматжилтын үнэлгээ, хуваарьт өдрүүдийн гүйцэтгэлийн тогтвортой байдал, гүйцэтгэх цаг ба газрын тогтвортой байдлыг нэгтгэсэн нийлмэл оноогоор илэрхийлэх загварыг боловсруулж, нотолгоо хангалтгүй үед оноо гаргахгүй байх дүрмийг тусгасан (3, 5-р бүлэг).
    \item Системийн туршилтын загварыг хэрэгжүүлж, 20 оролцогчтой урьдчилсан хэрэглэгчийн туршилт хийсэн (6-р бүлэг). Энэ зорилт хэсэгчлэн биелсэн: хэрэглээний тохиромжтой байдлын талаар хэрэглэгчийн үнэлгээ авсан боловч туршилтын хувилбарын хэмжилтийн дутагдлаас шалтгаалан зан үйлийн үзүүлэлтүүдийг дадал төлөвшлийн нотолгоо болгон ашиглах боломжгүй байв.
\end{enumerate}

\section*{Судалгааны асуултын хариулт}

\textbf{СА-1.} Онолын бүрэлдэхүүнүүдийг серверийн домэйн давхаргад өгөгдлийн сангаас хамааралгүй, тестлэгдэх дүрэм болгон хэрэгжүүлэх нь тохиромжтой болохыг энэ ажил харуулсан. Өдөөгч нөхцөл нь сануулга илгээх эсэхийг шийдэх дүрэм, доод зорилт нь гүйцэтгэлийг амжилттай гэж тооцох шалгуур, урамшуулал нь гүйцэтгэлийн бүртгэлтэй холбогдсон XP ба амжилтын тэмдэг, автоматжилт нь SRBAI үнэлгээ болон гүйцэтгэлийн эх үүсвэрийн ангиллаар тус тус илэрхийлэгдэв. Дүрэм бүр тодорхой босго, нотолгооны доод хэмжээтэй тул шийдвэрийг хэрэглэгчид тайлбарлах, нэгжийн тестээр шалгах боломжтой.

\textbf{СА-2.} Дадлын хүчийг ганц үзүүлэлтээр бус, хэрэглэгчийн өөрийн үнэлгээ (SRBAI) болон системийн бүртгэлээс гаргасан зан үйлийн прокси үзүүлэлтүүдийг жинлэн нэгтгэж илэрхийлэх аргыг санал болгов. Гүйцэтгэлийн тоо, дараалал нь дадал төлөвшсөнийг харуулахгүй тул тогтвортой байдлыг хуваарьт өдрүүдийн эзлэх хувиар, нөхцөл байдлын тогтвортой байдлыг гүйцэтгэх цаг ба газрын давтагдалтаар тооцсон. Энэ нийлмэл оноо болон түүний жингүүд нь онолд тулгуурласан эвристик бөгөөд дадлын автоматжилтыг хэмжих баталгаажсан хэмжүүр биш; түүний хүчинтэй байдлыг цаашдын судалгаагаар шалгах шаардлагатай.

\textbf{СА-3.} Урьдчилсан туршилтад оролцогчид дадал үүсгэх, гүйцэтгэл бүртгэх, сануулга авах, ахицаа харах үндсэн урсгалуудыг ойлгомжтой, хэрэгтэй гэж үнэлсэн бол дадал үүсгэх зааварчилгаа, дадлын хүчний онооны тайлбар, мэдэгдлийн тохиргоог сайжруулах хэрэгцээ тодорсон. Харин туршилтын хувилбарт өдрийн хил, гүйцэтгэлийн тогтвортой байдал, нөхцөл байдлын тогтвортой байдал, гүйцэтгэлийн эх үүсвэрийн хэмжилтэд дутагдал байсан тул дадал төлөвшилтэй холбоотой үзүүлэлтүүдийн хандлагын талаар найдвартай дүгнэлт гаргах боломжгүй байв. Эдгээр дутагдлыг эцсийн хувилбарт засварласан.

\section*{Судалгааны хувь нэмэр}

\begin{itemize}
    \item Дадал үүсэх онолын бүрэлдэхүүнүүдийг програм хангамжийн шаардлага, өгөгдлийн загвар, домэйн дүрэмтэй уялдуулсан зохиомжийн загвар.
    \item Цаг ба газрын өдөөгч нөхцөл, өдрийн хязгаар, гүйцэтгэлийн эх үүсвэрийн ангилал, дадлын хүчинд тулгуурласан сануулгын бууралтыг нэгтгэсэн, тайлбарлагдахуйц сануулгын логик.
    \item Нотолгоо хангалтгүй үед дүгнэлт гаргахгүй байх зарчимтай дадлын хүчний нийлмэл үнэлгээ.
    \item Монгол хэл дээрх, гар утас болон компьютерт ажиллах ажиллагаатай вэб програм болон түүний урьдчилсан хэрэглэгчийн үнэлгээ.
\end{itemize}

\section*{Хязгаарлалт}

Хэрэглэгчийн туршилт нь цөөн оролцогчтой, хүртээмжтэй түүвэрт тулгуурласан, хяналтын бүлэггүй, богино хугацааных байсан бөгөөд эцсийн хувилбараас өмнөх хувилбар дээр хийгдсэн. SRBAI-ийн монгол орчуулга хэмжилзүйн хувьд баталгаажаагүй. Нийлмэл онооны жин, сануулгын бууралтын босго нь онолд тулгуурласан эвристик утгууд юм. Байршлын өдөөгч нь хэрэглэгч програмыг нээсэн үед хуваалцсан сүүлийн байршилд тулгуурладаг тул програм хаалттай үед байршлыг хянахгүй. Иймээс систем дадал төлөвшилд нөлөөлдөг эсэхийг энэ ажил батлахгүй.

\section*{Цаашдын ажил}

\begin{itemize}
    \item Эцсийн хувилбар дээр хяналтын бүлэгтэй, олон долоо хоногийн туршилт хийж SRBAI, хуваарьт өдрүүдийн гүйцэтгэл, өөрийн санаачилгатай гүйцэтгэлийн хувийн өөрчлөлтийг хэмжих.
    \item SRBAI-ийн монгол орчуулгыг буцаан орчуулга, найдвартай байдлын шинжилгээгээр баталгаажуулах.
    \item Сануулгын бууралт өөрийн санаачилгатай гүйцэтгэлд хэрхэн нөлөөлж байгааг тусад нь үнэлэх.
    \item Нийлмэл онооны жин болон босгыг эмпирик өгөгдлөөр тохируулах.
    \item Хэрэглэгчийн саналд тулгуурлан дадал үүсгэх зааварчилгаа, онооны тайлбар, мэдэгдлийн тохиргоог сайжруулах.
\end{itemize}
